\documentclass[11pt]{ctexart}
\usepackage[a4paper,margin=1in]{geometry}
\usepackage{longtable,booktabs}
\title{Geant4-Agent Casebank Regression Report (EN)}
\author{Automated Pipeline}
\date{2026-03-08\_001313}
\begin{document}
\maketitle
\section{Unified Pass/Fail Table}
\begin{longtable}{p{0.06\linewidth}p{0.08\linewidth}p{0.07\linewidth}p{0.08\linewidth}p{0.07\linewidth}p{0.07\linewidth}p{0.07\linewidth}p{0.07\linewidth}p{0.09\linewidth}p{0.06\linewidth}p{0.21\linewidth}}
\toprule
ID & Domain & Style & Completeness & ExpInit & ActInit & ExpFinal & ActFinal & Missing(I->F) & Result & Fail Reason \\
\midrule
MTM01 & medical & fuzzy & missing & FAIL & FAIL & PASS & PASS & 12->0 & PASS & - \\
MTM02 & aerospace & precise & missing & FAIL & FAIL & PASS & PASS & 2->0 & PASS & - \\
MTM03 & industrial\_ndt & precise & missing & FAIL & FAIL & PASS & PASS & 12->0 & PASS & - \\
MTM04 & physics & precise & missing & FAIL & FAIL & PASS & PASS & 8->0 & PASS & - \\
MTM05 & nuclear\_engineering & fuzzy & missing & FAIL & FAIL & PASS & PASS & 6->0 & PASS & - \\
MTM06 & security\_screening & fuzzy & missing & FAIL & FAIL & PASS & PASS & 11->0 & PASS & - \\
MTM07 & semiconductor & precise & missing & FAIL & FAIL & PASS & PASS & 3->0 & PASS & - \\
MTM08 & education & fuzzy & missing & FAIL & FAIL & PASS & PASS & 11->0 & PASS & - \\
MTM09 & environmental\_radiation & fuzzy & missing & FAIL & FAIL & PASS & PASS & 3->0 & PASS & - \\
MTM10 & high\_energy\_physics & precise & missing & FAIL & FAIL & PASS & PASS & 10->0 & PASS & - \\
MTM11 & materials\_update & precise & full & PASS & PASS & PASS & PASS & 0->0 & PASS & - \\
MTM12 & overwrite\_reject & precise & full & PASS & PASS & PASS & PASS & 0->0 & PASS & - \\
MTM13 & medical\_detector & fuzzy & missing & FAIL & FAIL & PASS & PASS & 13->0 & PASS & - \\
MTM14 & industrial\_panel & precise & missing & FAIL & FAIL & PASS & PASS & 15->0 & PASS & - \\
MTM15 & shielding\_stack & precise & missing & FAIL & FAIL & PASS & PASS & 14->0 & PASS & - \\
MTM16 & nested\_target & precise & missing & FAIL & FAIL & PASS & PASS & 10->0 & PASS & - \\
MTM17 & coaxial\_shield & fuzzy & missing & FAIL & FAIL & PASS & PASS & 14->0 & PASS & - \\
MTM18 & boolean\_cavity & precise & missing & FAIL & FAIL & PASS & PASS & 14->0 & PASS & - \\
MTM19 & ring\_complete & precise & full & PASS & PASS & PASS & PASS & 0->0 & PASS & - \\
MTM20 & grid\_complete & precise & full & PASS & PASS & PASS & PASS & 0->0 & PASS & - \\
MTM21 & layered\_detector & fuzzy & missing & FAIL & FAIL & PASS & PASS & 14->0 & PASS & - \\
MTM22 & embedded\_sensor & precise & full & PASS & PASS & PASS & PASS & 0->0 & PASS & - \\
MTM23 & layered\_shield\_complete & precise & full & PASS & PASS & PASS & PASS & 0->0 & PASS & - \\
MTM24 & boolean\_complete & precise & full & PASS & PASS & PASS & PASS & 0->0 & PASS & - \\
\bottomrule
\end{longtable}
\section{Baseline Comparison}
\begin{itemize}
\item pass\_rate=1.000
\item initial\_alignment\_rate=1.000
\item final\_alignment\_rate=1.000
\item avg\_turns\_per\_case=2.458
\item missing\_reduction\_rate=1.000
\item confirm\_apply\_success\_rate=1.000
\item rollback\_after\_confirm\_rate=0.000
\item internal\_leak\_turn\_rate=0.034
\item repeated\_assistant\_case\_rate=0.000
\item stale\_confirm\_case\_rate=0.000
\item delta\_pass\_rate=+0.050
\item regressions=none
\item improvements=none
\end{itemize}
\appendix
\section{Appendix: Full Dialogues and Naturalized Internal Chains}
\subsection{Case MTM01 (medical)}
\textbf{Dialogue:}
\begin{itemize}
\item User: 我先要一个用于医学成像的初步辐照场景，细节还没定，你先告诉我还缺什么。
\item Assistant: 目前我们已经确定了使用盒体作为几何结构，接下来还需要补充以下信息：

1.  **几何**：盒体在 x、y、z 三个方向的具体尺寸。
2.  **材料**：需要为这个盒体体积指定对应的材料。
3.  **源**：需要确定放射源的类型（例如伽马源）、粒子种类和能量。

请随时告诉我这些细节，我们就能继续推进配置了。
\item User: 采用 300 mm x 300 mm x 300 mm 的水箱，材料 G4\_WATER；源用束流 gamma，能量 6 MeV，位置 (0,0,-400) mm，方向 +z；物理列表 QBBC；输出 root。
\item Assistant: 我注意到这可能会覆盖已经确认的设置。具体来说：几何结构类型将从“未知”改为“单盒体”；盒体 x 方向尺寸将从“未设置”改为“300.0”。如果您确认应用这些修改，请直接回复“确认”；如果您希望保持原有的值，请直接回复“保持原值”。
\item User: 确认
\item Assistant: 配置已完成。当前已确认的设置有：盒体在 x、y、z 三个方向的尺寸以及几何结构类型。
\end{itemize}
\textbf{Internal Chain (Natural Language):}
\begin{itemize}
\item Turn 1: action=summarize\_progress, phase=geometry, complete=False.
\item Missing fields: geometry.params.module\_x, geometry.params.module\_y, geometry.params.module\_z, materials.volume\_material\_map, source.type, source.particle, source.energy, source.position, source.direction, physics.physics\_list, output.format, output.path.
\item LLM used=True; fallback=none; backend=llm\_semantic\_frame+runtime\_semantic.
\item Turn 2: action=confirm\_overwrite, phase=geometry, complete=False.
\item Missing fields: geometry.params.module\_x, geometry.params.module\_y, geometry.params.module\_z.
\item LLM used=True; fallback=none; backend=llm\_slot\_frame+runtime\_semantic.
\item Turn 3: action=finalize, phase=materials, complete=True.
\item Missing fields: none.
\item LLM used=True; fallback=none; backend=llm\_slot\_frame+runtime\_semantic+confirmed\_pending\_overwrite.
\end{itemize}
\subsection{Case MTM02 (aerospace)}
\textbf{Dialogue:}
\begin{itemize}
\item User: I need a first-pass proton irradiation setup for a silicon device; keep it simple for now.
\item Assistant: Great! I've updated your configuration based on your request for a simple proton irradiation setup for a silicon device.

Here's what's been completed:
*   **Geometry:** A box structure with specified X and Y dimensions.
*   **Materials:** Silicon has been selected and assigned to the volume.
*   **Source:** A proton source with a defined energy is configured.
*   **Physics:** An appropriate physics list has been set.

To finalize the setup, I still need:
*   **Output:** The desired output format and a file path to save the results.
\item User: Use a silicon box 50 mm x 50 mm x 5 mm, beam source proton 100 MeV from (0,0,-100) mm along +z, physics FTFP\_BERT, output json.
\item Assistant: Your request would change a previously confirmed geometry setting.

Pending change:
*   Box size along X: 1000.0 mm → 50.0 mm
*   Box size along Y: 1000.0 mm → 50.0 mm

Please reply **'confirm'** to apply these changes, or **'keep existing'** to keep the current values.
\item User: confirm
\item Assistant: Configuration complete. The following items are confirmed: source energy, box size along Z, box size along Y, and box size along X.
\end{itemize}
\textbf{Internal Chain (Natural Language):}
\begin{itemize}
\item Turn 1: action=summarize\_progress, phase=materials, complete=False.
\item Missing fields: output.format, output.path.
\item LLM used=True; fallback=none; backend=llm\_slot\_frame+runtime\_semantic.
\item Turn 2: action=confirm\_overwrite, phase=source, complete=False.
\item Missing fields: none.
\item LLM used=True; fallback=none; backend=llm\_slot\_frame+runtime\_semantic.
\item Turn 3: action=finalize, phase=physics, complete=True.
\item Missing fields: none.
\item LLM used=True; fallback=none; backend=llm\_semantic\_frame+runtime\_semantic+confirmed\_pending\_overwrite.
\end{itemize}
\subsection{Case MTM03 (industrial\_ndt)}
\textbf{Dialogue:}
\begin{itemize}
\item User: 我要做一个 gamma 穿透铁靶的小测试，但我还没把完整参数说完。
\item Assistant: 目前我们已确认了材料（铁）和粒子类型（伽马射线）。接下来还需要您提供几何结构的具体类型和尺寸，以及体积与材料的绑定方式。请继续补充这些信息，我们就能继续推进配置了。
\item User: 请设为半径 40 mm、半长 80 mm 的铁圆柱，材料 G4\_Fe；点源 gamma，能量 2 MeV，位置 (0,0,-200) mm，方向 +z；物理列表 Shielding；输出 csv。
\item Assistant: 配置已完成。当前已确认的设置有：盒体在 x、y、z 方向的尺寸，以及体积与材料的绑定关系。
\end{itemize}
\textbf{Internal Chain (Natural Language):}
\begin{itemize}
\item Turn 1: action=summarize\_progress, phase=geometry, complete=False.
\item Missing fields: geometry.structure, geometry.params.module\_x, geometry.params.module\_y, geometry.params.module\_z, materials.volume\_material\_map, source.type, source.energy, source.position, source.direction, physics.physics\_list, output.format, output.path.
\item LLM used=True; fallback=none; backend=llm\_slot\_frame+runtime\_semantic.
\item Turn 2: action=finalize, phase=materials, complete=True.
\item Missing fields: none.
\item LLM used=True; fallback=none; backend=llm\_slot\_frame+runtime\_semantic.
\end{itemize}
\subsection{Case MTM04 (physics)}
\textbf{Dialogue:}
\begin{itemize}
\item User: Gamma point source 1 MeV at (0,0,-100) toward +z; please build the rest.
\item Assistant: Geometry: still needs geometry structure, box size along x. Materials: still needs volume-to-material binding. Source: updated source type, particle type, source energy. Physics: still needs physics list. Output: still needs output format, output path.
\item User: Set copper box 1 m x 1 m x 1 m, physics FTFP\_BERT, output hdf5.
\item Assistant: Your configuration is complete. The following items are confirmed: volume-to-material binding, output path, output format, and physics list.
\end{itemize}
\textbf{Internal Chain (Natural Language):}
\begin{itemize}
\item Turn 1: action=summarize\_progress, phase=geometry, complete=False.
\item Missing fields: geometry.structure, geometry.params.module\_x, geometry.params.module\_y, geometry.params.module\_z, materials.volume\_material\_map, physics.physics\_list, output.format, output.path.
\item LLM used=True; fallback=none; backend=llm\_slot\_frame+runtime\_semantic.
\item Turn 2: action=finalize, phase=materials, complete=True.
\item Missing fields: none.
\item LLM used=True; fallback=none; backend=llm\_slot\_frame+runtime\_semantic.
\end{itemize}
\subsection{Case MTM05 (nuclear\_engineering)}
\textbf{Dialogue:}
\begin{itemize}
\item User: 我想做一个中子屏蔽的概念验证，现在先给一个缺参开头。
\item Assistant: 好的，我们已同步了以下设置：
*   **几何**：结构类型
*   **材料**：材料选择及绑定
*   **源**：粒子类型和源类型
*   **物理**：物理列表
*   **输出**：格式和路径

接下来，还需要您补充几个关键参数：
*   几何部分的 **盒体尺寸**（x, y, z 方向）
*   源部分的 **能量** 和 **位置**

请提供这些信息，我们就能继续推进配置了。
\item User: 使用 200 mm x 200 mm x 200 mm 的铅块，材料 G4\_Pb；点源 neutron，能量 14 MeV，位置 (0,0,-50) mm，方向 +z；物理列表 QGSP\_BERT\_HP；输出 xml。
\item Assistant: 我检测到这会覆盖已经确认的值。待修改内容：几何结构类型：unknown -> single\_box；已选材料：['G4\_Cu'] -> ['G4\_Pb']。如果要应用这些修改，请直接回复“确认”；如果要保留原值，请直接说“保持原值”。
\item User: 确认
\item Assistant: 配置已完成。当前已确认的参数包括：体积与材料绑定、输出路径、盒体 z 方向尺寸、盒体 y 方向尺寸。
\end{itemize}
\textbf{Internal Chain (Natural Language):}
\begin{itemize}
\item Turn 1: action=summarize\_progress, phase=geometry, complete=False.
\item Missing fields: geometry.params.module\_x, geometry.params.module\_y, geometry.params.module\_z, source.energy, source.position, source.direction.
\item LLM used=True; fallback=none; backend=llm\_slot\_frame+runtime\_semantic.
\item Turn 2: action=confirm\_overwrite, phase=geometry, complete=False.
\item Missing fields: geometry.params.module\_x, geometry.params.module\_y, geometry.params.module\_z.
\item LLM used=True; fallback=none; backend=llm\_slot\_frame+runtime\_semantic.
\item Turn 3: action=finalize, phase=materials, complete=True.
\item Missing fields: none.
\item LLM used=True; fallback=none; backend=llm\_slot\_frame+runtime\_semantic+confirmed\_pending\_overwrite.
\end{itemize}
\subsection{Case MTM06 (security\_screening)}
\textbf{Dialogue:}
\begin{itemize}
\item User: Need a quick cargo-screening toy setup with forward photons, but I have not fixed the target shape, material, source energy, or output format yet.
\item Assistant: Great. Here's what we have so far:

**Completed:**
*   **Geometry:** The structure is set.
*   **Source:** The particle type (photon) is confirmed.

**Still Needed:**
*   **Geometry:** The box dimensions (size along X, Y, and Z).
*   **Materials:** The material to fill the target volume.
*   **Source:** The specific source type and its energy.
*   **Physics:** A physics list.
*   **Output:** The output format and file path.
\item User: Use steel cylinder target radius 60 mm half-length 120 mm, beam gamma 3 MeV from (0,0,-500) mm to +z, physics Shielding, output root.
\item Assistant: This would overwrite an already confirmed value. Pending change: geometry structure: unknown -> single\_tubs; box size along x: None -> 1000.0. Reply 'confirm' to apply it, or say 'keep existing' to keep the current value.
\item User: confirm
\item Assistant: Configuration complete. Confirmed items: volume-to-material binding, half length, radius, and box size along Z.
\end{itemize}
\textbf{Internal Chain (Natural Language):}
\begin{itemize}
\item Turn 1: action=summarize\_progress, phase=geometry, complete=False.
\item Missing fields: geometry.params.module\_x, geometry.params.module\_y, geometry.params.module\_z, materials.volume\_material\_map, source.type, source.energy, source.position, source.direction, physics.physics\_list, output.format, output.path.
\item LLM used=True; fallback=none; backend=llm\_slot\_frame+runtime\_semantic.
\item Turn 2: action=confirm\_overwrite, phase=geometry, complete=False.
\item Missing fields: geometry.params.module\_x, geometry.params.module\_y, geometry.params.module\_z.
\item LLM used=True; fallback=none; backend=llm\_slot\_frame+runtime\_semantic.
\item Turn 3: action=finalize, phase=materials, complete=True.
\item Missing fields: none.
\item LLM used=True; fallback=none; backend=llm\_semantic\_frame+runtime\_semantic+confirmed\_pending\_overwrite.
\end{itemize}
\subsection{Case MTM07 (semiconductor)}
\textbf{Dialogue:}
\begin{itemize}
\item User: 我想搭一个硅薄片的中子照射测试，但先只给一个简述。
\item Assistant: 目前我们已经同步了以下配置：

- **几何**：盒体结构及其尺寸
- **材料**：已选定硅材料并绑定到体积
- **源**：中子源的类型、能量和位置已设定

接下来还需要您补充以下信息才能完成配置：
- 物理列表
- 输出格式
- 输出路径
\item User: 使用 20 mm x 20 mm x 1 mm 的硅片，材料 G4\_Si；各向同性 neutron，能量 2 MeV，位置 (0,0,0) mm；物理列表 QGSP\_BERT\_HP；输出 json。
\item Assistant: 我注意到这次修改会覆盖已确认的配置。具体变动如下：

- 盒体 x 方向尺寸：从 1000.0 改为 20.0
- 盒体 y 方向尺寸：从 1000.0 改为 20.0

请确认您的选择：
- 如果同意应用这些修改，请直接回复“确认”。
- 如果希望保留原来的数值，请直接回复“保持原值”。
\item User: 确认
\item Assistant: 配置已完成。当前已确认的设置有：源位置、源能量、源类型以及盒体 z 方向尺寸。
\end{itemize}
\textbf{Internal Chain (Natural Language):}
\begin{itemize}
\item Turn 1: action=summarize\_progress, phase=materials, complete=False.
\item Missing fields: physics.physics\_list, output.format, output.path.
\item LLM used=True; fallback=none; backend=llm\_slot\_frame+runtime\_semantic.
\item Turn 2: action=confirm\_overwrite, phase=source, complete=False.
\item Missing fields: none.
\item LLM used=True; fallback=none; backend=llm\_slot\_frame+runtime\_semantic.
\item Turn 3: action=finalize, phase=physics, complete=True.
\item Missing fields: none.
\item LLM used=True; fallback=none; backend=llm\_semantic\_frame+runtime\_semantic+confirmed\_pending\_overwrite.
\end{itemize}
\subsection{Case MTM08 (education)}
\textbf{Dialogue:}
\begin{itemize}
\item User: For classroom demo, I want gamma attenuation in copper; you can ask only essential missing parts.
\item Assistant: Here's a summary of your setup so far:

**Completed:**
*   **Materials:** Copper is selected.
*   **Source:** Gamma rays are set as the particle type.
*   **Physics:** A physics list has been configured.

**Still needed:**
*   **Geometry:** The structure type and the box size along the X, Y, and Z axes.
*   **Materials:** You'll need to assign the copper material to a specific volume.
*   **Source:** The source type, its energy, and its position.
*   **Output:** The format and file path for your results.
\item User: Okay, copper box 100 mm x 100 mm x 100 mm, point gamma 1 MeV at (0,0,-200) mm, direction +z, physics FTFP\_BERT, output csv.
\item Assistant: This would change a setting you've already confirmed.

**Pending change:**  
Physics list: `Shielding` → `FTFP\_BERT`

Reply **"confirm"** to apply this change, or **"keep existing"** to keep the current value.
\item User: confirm
\item Assistant: Your configuration is complete.

The following items are confirmed: physics list, output path, output format, and source direction.
\end{itemize}
\textbf{Internal Chain (Natural Language):}
\begin{itemize}
\item Turn 1: action=summarize\_progress, phase=geometry, complete=False.
\item Missing fields: geometry.structure, geometry.params.module\_x, geometry.params.module\_y, geometry.params.module\_z, materials.volume\_material\_map, source.type, source.energy, source.position, source.direction, output.format, output.path.
\item LLM used=True; fallback=none; backend=llm\_slot\_frame+runtime\_semantic.
\item Turn 2: action=confirm\_overwrite, phase=materials, complete=False.
\item Missing fields: none.
\item LLM used=True; fallback=none; backend=llm\_slot\_frame+runtime\_semantic.
\item Turn 3: action=finalize, phase=source, complete=True.
\item Missing fields: none.
\item LLM used=True; fallback=none; backend=llm\_semantic\_frame+runtime\_semantic+confirmed\_pending\_overwrite.
\end{itemize}
\subsection{Case MTM09 (environmental\_radiation)}
\textbf{Dialogue:}
\begin{itemize}
\item User: 我想先搭一个空气体积中的低能 gamma 传播场景，细节稍后补。
\item Assistant: 好的，目前配置进度如下：

**已完成设置：**
*   **几何**：已设定为长方体，并确定了X、Y方向的尺寸。
*   **材料**：已选定材料，并完成了与体积的绑定。
*   **源**：已设定粒子类型、能量和位置。

**仍需补充的信息：**
*   **物理**：请选择物理列表。
*   **输出**：请指定输出格式和输出文件的保存路径。

请提供上述信息以继续完成配置。
\item User: 设为 500 mm x 500 mm x 500 mm 的空气盒，材料 G4\_AIR；点源 gamma，能量 0.5 MeV，位置 (0,0,-300) mm，方向 +z；物理列表 QBBC；输出 root。
\item Assistant: 我注意到这将会覆盖已经确认的参数。具体修改如下：

*   盒体 x 方向尺寸：1000.0 → 500.0
*   盒体 y 方向尺寸：1000.0 → 500.0

请确认您的选择：如果应用这些修改，请回复“确认”；如果想保留原值，请回复“保持原值”。
\item User: 确认
\item Assistant: 配置已完成。当前已确认的参数包括：源位置、源能量、盒体 z 方向尺寸、盒体 y 方向尺寸和盒体 x 方向尺寸。
\end{itemize}
\textbf{Internal Chain (Natural Language):}
\begin{itemize}
\item Turn 1: action=summarize\_progress, phase=materials, complete=False.
\item Missing fields: physics.physics\_list, output.format, output.path.
\item LLM used=True; fallback=none; backend=llm\_slot\_frame+runtime\_semantic.
\item Turn 2: action=confirm\_overwrite, phase=source, complete=False.
\item Missing fields: none.
\item LLM used=True; fallback=none; backend=llm\_slot\_frame+runtime\_semantic.
\item Turn 3: action=finalize, phase=physics, complete=True.
\item Missing fields: none.
\item LLM used=True; fallback=none; backend=llm\_semantic\_frame+runtime\_semantic+confirmed\_pending\_overwrite.
\end{itemize}
\subsection{Case MTM10 (high\_energy\_physics)}
\textbf{Dialogue:}
\begin{itemize}
\item User: Set geometry as a sphere target; I have not decided radius or source details yet.
\item Assistant: Here's where we stand:

**Geometry:** You've set the target to a sphere. We still need the sphere's radius.

**Materials:** We need to assign a material to your volume.

**Source:** We need the source type, particle type, and particle energy.

**Physics \& Output:** We also need the physics list, output format, and file path.

What would you like to provide next?
\item User: Sphere target radius 80 mm, lead material, electron point source 5 MeV at (0,0,-200) mm toward +z, physics FTFP\_BERT, output hdf5.
\item Assistant: Your configuration is complete. The following items are confirmed: volume-to-material binding, output path, output format, and physics list.
\end{itemize}
\textbf{Internal Chain (Natural Language):}
\begin{itemize}
\item Turn 1: action=summarize\_progress, phase=geometry, complete=False.
\item Missing fields: geometry.params.child\_rmax, materials.volume\_material\_map, source.type, source.particle, source.energy, source.position, source.direction, physics.physics\_list, output.format, output.path.
\item LLM used=True; fallback=none; backend=llm\_slot\_frame+runtime\_semantic.
\item Turn 2: action=finalize, phase=materials, complete=True.
\item Missing fields: none.
\item LLM used=True; fallback=none; backend=llm\_slot\_frame+runtime\_semantic.
\end{itemize}
\subsection{Case MTM11 (materials\_update)}
\textbf{Dialogue:}
\begin{itemize}
\item User: Create a complete setup: 100 mm x 100 mm x 100 mm copper box, point gamma 1 MeV at (0,0,-200) mm toward +z, physics QBBC, output json.
\item Assistant: Your Geant4 configuration is complete. The following items have been confirmed: the assignment of materials to volumes, the output file path, and the JSON output format.
\item User: Change the target material to G4\_Al and change the output format to root.
\item Assistant: You're about to overwrite two confirmed settings:
- Selected materials: from copper (G4\_Cu) to aluminum (G4\_Al)
- Output format: from JSON to ROOT

Please reply **'confirm'** to apply these changes, or **'keep existing'** to keep the current values.
\item User: confirm
\item Assistant: Your configuration is complete. The following items are confirmed: the assignment of materials to volumes, the output file path, the ROOT output format, and the selected materials (including the change to aluminum).
\end{itemize}
\textbf{Internal Chain (Natural Language):}
\begin{itemize}
\item Turn 1: action=finalize, phase=materials, complete=True.
\item Missing fields: none.
\item LLM used=True; fallback=none; backend=llm\_slot\_frame+runtime\_semantic.
\item Turn 2: action=confirm\_overwrite, phase=source, complete=False.
\item Missing fields: none.
\item LLM used=True; fallback=none; backend=llm\_slot\_frame+runtime\_semantic.
\item Turn 3: action=finalize, phase=physics, complete=True.
\item Missing fields: none.
\item LLM used=True; fallback=none; backend=llm\_semantic\_frame+runtime\_semantic+confirmed\_pending\_overwrite.
\end{itemize}
\subsection{Case MTM12 (overwrite\_reject)}
\textbf{Dialogue:}
\begin{itemize}
\item User: 请创建一个完整配置：100 mm x 100 mm x 100 mm 的铜盒，材料 G4\_Cu；点源 gamma，能量 1 MeV，位置 (0,0,-200) mm，方向 +z；物理列表 QBBC；输出 json。
\item Assistant: 配置已完成。当前已确认的设置包括：体积与材料绑定、输出路径和输出格式。
\item User: 把材料改成 G4\_Al，并把输出格式改成 root。
\item Assistant: 我注意到这次修改会覆盖已经确认的设置。具体来说：

- 已选材料：从 `G4\_Cu` 改为 `G4\_Al`
- 输出格式：从 `json` 改为 `root`

如果你确定要应用这些修改，请直接回复“确认”；如果想保留原来的值，请直接说“保持原值”。
\item User: 保持原值
\item Assistant: 好的，已保留已选材料和输出格式的原有值，这次修改没有应用。
\end{itemize}
\textbf{Internal Chain (Natural Language):}
\begin{itemize}
\item Turn 1: action=finalize, phase=materials, complete=True.
\item Missing fields: none.
\item LLM used=True; fallback=none; backend=llm\_slot\_frame+runtime\_semantic.
\item Turn 2: action=confirm\_overwrite, phase=source, complete=False.
\item Missing fields: geometry.ask.boolean.solid\_a\_size, geometry.ask.boolean.solid\_b\_size.
\item LLM used=True; fallback=none; backend=llm\_slot\_frame+runtime\_semantic.
\item Turn 3: action=reject\_overwrite, phase=physics, complete=True.
\item Missing fields: none.
\item LLM used=True; fallback=none; backend=llm\_semantic\_frame+runtime\_semantic+rejected\_pending\_overwrite.
\end{itemize}
\subsection{Case MTM13 (medical\_detector)}
\textbf{Dialogue:}
\begin{itemize}
\item User: I want a ring of detector modules around a small source region, but I have not fixed the ring radius or module size yet.
\item Assistant: Your geometry structure is now set to a ring of detector modules. To continue, I still need:

*   **Geometry:** The size of each ring module and how many modules are in the ring.
*   **Materials:** To assign materials to each volume.
*   **Source:** The type of source, the particle type, and its energy.
*   **Physics:** The physics list to use.
*   **Output:** The format for results and where to save them.
\item User: Use a ring of 12 modules, each 8 mm x 10 mm x 2 mm, radius 40 mm, clearance 1 mm, material G4\_WATER; point gamma source 0.662 MeV at (0,0,0) mm; physics QBBC; output json.
\item Assistant: This would overwrite an already confirmed geometry. The pending change is:

**From:**
*   Module: Box (7.83 mm x 4.72 mm x 6.86 mm)
*   Ring: 4 modules, radius 67.49 mm, clearance 0.19 mm

**To:**
*   Module: Box (8.0 mm x 10.0 mm x 2.0 mm)
*   Ring: 12 modules, radius 40.0 mm, clearance 1.0 mm

Reply **'confirm'** to apply this change, or **'keep existing'** to keep the current values.
\item User: confirm
\item Assistant: Configuration complete. The ring geometry is now set with your confirmed values for clearance, radius, module count, and module size.
\end{itemize}
\textbf{Internal Chain (Natural Language):}
\begin{itemize}
\item Turn 1: action=summarize\_progress, phase=materials, complete=False.
\item Missing fields: materials.volume\_material\_map, source.type, source.particle, source.energy, source.position, source.direction, physics.physics\_list, output.format, output.path, geometry.ask.ring.module\_size, geometry.ask.ring.count, geometry.ask.ring.radius, geometry.ask.ring.clearance.
\item LLM used=True; fallback=none; backend=llm\_slot\_frame+runtime\_semantic.
\item Turn 2: action=confirm\_overwrite, phase=source, complete=False.
\item Missing fields: none.
\item LLM used=True; fallback=none; backend=llm\_slot\_frame+runtime\_semantic.
\item Turn 3: action=finalize, phase=physics, complete=True.
\item Missing fields: none.
\item LLM used=True; fallback=none; backend=llm\_semantic\_frame+runtime\_semantic+confirmed\_pending\_overwrite.
\end{itemize}
\subsection{Case MTM14 (industrial\_panel)}
\textbf{Dialogue:}
\begin{itemize}
\item User: 我想做一个二维阵列探测板，但 pitch 和材料还没定。
\item Assistant: 好的，我们已确认了您要创建二维阵列探测板。目前几何部分已明确结构类型，但还需要您提供阵列模块的尺寸和X方向的阵列数量。

接下来，为了完成配置，我们还需要您补充以下信息：
*   **材料**：每个体积对应的材料。
*   **源**：源的类型、粒子种类和能量。
*   **物理**：需要使用的物理过程列表。
*   **输出**：希望输出的文件格式和保存路径。
\item User: 设成 4 x 3 的盒体阵列，每个模块 10 mm x 10 mm x 2 mm，pitch\_x 12 mm，pitch\_y 14 mm，间隙 1 mm，材料 G4\_Si；点源 gamma 0.511 MeV，位置 (0,0,-80) mm，方向 +z；物理列表 QBBC；输出 root。
\item Assistant: 检测到您当前的设置会覆盖已确认的几何参数。具体来说，阵列模块的尺寸、X/Y方向的间距以及阵列间隙都将被更新为新的数值。

如果您确认应用这些修改，请直接回复“确认”；如果您希望保留原来的数值，请直接说“保持原值”。
\item User: 确认
\item Assistant: 配置已完成！目前几何部分的以下参数已确认：y 方向阵列数量、x 方向阵列数量、阵列间隙、y 方向 pitch。
\end{itemize}
\textbf{Internal Chain (Natural Language):}
\begin{itemize}
\item Turn 1: action=summarize\_progress, phase=materials, complete=False.
\item Missing fields: materials.volume\_material\_map, source.type, source.particle, source.energy, source.position, source.direction, physics.physics\_list, output.format, output.path, geometry.ask.grid.module\_size, geometry.ask.grid.count\_x, geometry.ask.grid.count\_y, geometry.ask.grid.pitch\_x, geometry.ask.grid.pitch\_y, geometry.ask.grid.clearance.
\item LLM used=True; fallback=none; backend=llm\_semantic\_frame+runtime\_semantic.
\item Turn 2: action=confirm\_overwrite, phase=source, complete=False.
\item Missing fields: geometry.ask.grid.count\_x, geometry.ask.grid.count\_y.
\item LLM used=True; fallback=none; backend=llm\_slot\_frame+runtime\_semantic.
\item Turn 3: action=finalize, phase=physics, complete=True.
\item Missing fields: none.
\item LLM used=True; fallback=none; backend=llm\_semantic\_frame+runtime\_semantic+confirmed\_pending\_overwrite.
\end{itemize}
\subsection{Case MTM15 (shielding\_stack)}
\textbf{Dialogue:}
\begin{itemize}
\item User: Build a stacked absorber sandwich along z; I will provide the layer sizes next.
\item Assistant: Geometry updated: you've set a stacked absorber sandwich structure along the z-axis. Next, I'll need the stack layer footprint and thicknesses.

Still to configure:
- Materials: volume-to-material binding.
- Source: type, particle type, and energy.
- Physics: physics list.
- Output: format and file path.
\item User: Use stacked layers with footprint 20 mm x 20 mm, thicknesses 2 mm, 4 mm, 6 mm, layer clearance 0.5 mm, outer box 30 mm x 30 mm x 20 mm, container clearance 1 mm, material G4\_Al; beam gamma 1.25 MeV from (0,0,-120) mm along +z; physics Shielding; output csv.
\item Assistant: Your new geometry settings would overwrite the current confirmed geometry.

**Current values:**
- Stack footprint: \textasciitilde{}16.33 mm x \textasciitilde{}10.28 mm
- Layer thicknesses: \textasciitilde{}6.68 mm, \textasciitilde{}1.19 mm, \textasciitilde{}5.59 mm
- Layer clearance: \textasciitilde{}0.73 mm
- Outer box: \textasciitilde{}18.07 mm x \textasciitilde{}45.52 mm x \textasciitilde{}16.66 mm
- Container clearance: \textasciitilde{}0.87 mm

**New values:**
- Stack footprint: 20.0 mm x 20.0 mm
- Layer thicknesses: 2.0 mm, 4.0 mm, 6.0 mm
- Layer clearance: 0.5 mm
- Outer box: 30.0 mm x 30.0 mm x 20.0 mm
- Container clearance: 1.0 mm

Reply **'confirm'** to apply these new values, or **'keep existing'** to keep the current geometry.
\item User: confirm
\item Assistant: Your configuration is complete. The following geometry items are confirmed: stack container clearance, stack container size, stack layer clearance, and stack layer thicknesses.
\end{itemize}
\textbf{Internal Chain (Natural Language):}
\begin{itemize}
\item Turn 1: action=summarize\_progress, phase=materials, complete=False.
\item Missing fields: materials.volume\_material\_map, source.type, source.particle, source.energy, source.position, source.direction, physics.physics\_list, output.format, output.path, geometry.ask.stack.footprint, geometry.ask.stack.thicknesses, geometry.ask.stack.layer\_clearance, geometry.ask.stack.parent\_size, geometry.ask.stack.nest\_clearance.
\item LLM used=True; fallback=none; backend=llm\_slot\_frame+runtime\_semantic.
\item Turn 2: action=confirm\_overwrite, phase=source, complete=False.
\item Missing fields: none.
\item LLM used=True; fallback=none; backend=llm\_slot\_frame+runtime\_semantic.
\item Turn 3: action=finalize, phase=physics, complete=True.
\item Missing fields: none.
\item LLM used=True; fallback=none; backend=llm\_semantic\_frame+runtime\_semantic+confirmed\_pending\_overwrite.
\end{itemize}
\subsection{Case MTM16 (nested\_target)}
\textbf{Dialogue:}
\begin{itemize}
\item User: 我想在一个盒子里嵌一个圆柱靶，先别急着补全。
\item Assistant: 目前我们已经确认了：几何结构为“盒子内嵌圆柱”，以及源的位置、方向和类型。

接下来还需要补充以下信息才能完成配置：
1. 几何部分：外层盒子的具体尺寸，以及内嵌圆柱的半径。
2. 材料部分：为各个体积指定材料。
3. 源部分：粒子类型和能量。
4. 物理部分：选择物理列表。
5. 输出部分：输出格式和文件路径。

请随时告诉我这些信息。
\item User: 外盒 120 mm x 120 mm x 120 mm，内嵌铜圆柱半径 20 mm、半长 40 mm，嵌套间隙 2 mm；材料 G4\_Cu；束流 gamma 1 MeV，位置 (0,0,-150) mm，方向 +z；物理列表 QBBC；输出 json。
\item Assistant: 检测到以下已确认的设置将被覆盖：
- 源类型：从“点源”改为“束流”
- 源位置：从 [0, 0, -100] mm 改为 [0, 0, -150] mm

请确认是否应用这些修改：
- 若同意覆盖，请回复“确认”。
- 若希望保留原值，请回复“保持原值”。
\item User: 确认
\item Assistant: 配置已完成。当前已确认的设置包括：源位置、源类型、嵌套间隙和内嵌圆柱半长。
\end{itemize}
\textbf{Internal Chain (Natural Language):}
\begin{itemize}
\item Turn 1: action=summarize\_progress, phase=materials, complete=False.
\item Missing fields: materials.volume\_material\_map, source.particle, source.energy, physics.physics\_list, output.format, output.path, geometry.ask.nest.parent\_size, geometry.ask.nest.child\_radius, geometry.ask.nest.child\_half\_length, geometry.ask.nest.clearance.
\item LLM used=True; fallback=none; backend=llm\_slot\_frame+runtime\_semantic.
\item Turn 2: action=confirm\_overwrite, phase=source, complete=False.
\item Missing fields: none.
\item LLM used=True; fallback=none; backend=llm\_slot\_frame+runtime\_semantic.
\item Turn 3: action=finalize, phase=physics, complete=True.
\item Missing fields: none.
\item LLM used=True; fallback=none; backend=llm\_semantic\_frame+runtime\_semantic+confirmed\_pending\_overwrite.
\end{itemize}
\subsection{Case MTM17 (coaxial\_shield)}
\textbf{Dialogue:}
\begin{itemize}
\item User: I need a shell-style shielding toy model, but the shell thicknesses are still open.
\item Assistant: Here's what we have so far:

**Geometry:** The structure is set to a shell. We still need the shell's inner radius, thicknesses, and half-length.
**Materials:** We need to assign materials to the volumes.
**Source:** We need the source type, particle type, energy, and position.
**Physics:** The physics list is configured.
**Output:** We need the output format and file path.
\item User: Use a shell with inner radius 15 mm, thicknesses 5 mm and 8 mm, half length 40 mm, nested core radius 10 mm, nested core half length 30 mm, clearance 1 mm, material G4\_Pb; isotropic neutron source 2 MeV at (0,0,0) mm; physics Shielding; output hdf5.
\item Assistant: This would overwrite a previously confirmed geometry.

**Pending change:**
*   **Shell:** Inner radius from 7.15 mm to 15.0 mm, thicknesses from [1.18, 3.43, 0.86] mm to [5.0, 8.0, 0.83] mm, half-length from 22.36 mm to 40.0 mm.
*   **Nested Core:** Radius from 11.60 mm to 10.0 mm, half-length from 2.68 mm to 30.0 mm.
*   **Clearance:** From 1.01 mm to 1.0 mm.

Reply **'confirm'** to apply these new values, or **'keep existing'** to keep the current geometry.
\item User: confirm
\item Assistant: Your configuration is complete. The shell clearance, nested core dimensions, and shell half length have been confirmed.
\end{itemize}
\textbf{Internal Chain (Natural Language):}
\begin{itemize}
\item Turn 1: action=summarize\_progress, phase=materials, complete=False.
\item Missing fields: materials.volume\_material\_map, source.type, source.particle, source.energy, source.position, source.direction, output.format, output.path, geometry.ask.shell.inner\_radius, geometry.ask.shell.thicknesses, geometry.ask.shell.half\_length, geometry.ask.shell.child\_radius, geometry.ask.shell.child\_half\_length, geometry.ask.shell.clearance.
\item LLM used=True; fallback=none; backend=llm\_semantic\_frame+runtime\_semantic.
\item Turn 2: action=confirm\_overwrite, phase=source, complete=False.
\item Missing fields: none.
\item LLM used=True; fallback=none; backend=llm\_slot\_frame+runtime\_semantic.
\item Turn 3: action=finalize, phase=physics, complete=True.
\item Missing fields: none.
\item LLM used=True; fallback=none; backend=llm\_semantic\_frame+runtime\_semantic+confirmed\_pending\_overwrite.
\end{itemize}
\subsection{Case MTM18 (boolean\_cavity)}
\textbf{Dialogue:}
\begin{itemize}
\item User: 我想做一个挖空的盒体几何，但先只给你意图。
\item Assistant: 好的，我们已同步了您关于几何结构、源位置、方向和类型的设置。接下来，还需要您补充挖空盒体（Boolean Solid A 和 B）的具体尺寸信息，以及材料、物理过程等配置。
\item User: 做减法布尔体：A 为 40 mm x 40 mm x 40 mm 盒体，B 为 20 mm x 20 mm x 20 mm 盒体，从 A 中减去 B；材料 G4\_Al；点源 gamma 0.3 MeV，位置 (0,0,-60) mm，方向 +z；物理列表 QBBC；输出 xml。
\item Assistant: 我注意到这可能会覆盖已经确认的设置。具体来说：

- **源位置**：将从 (0,0,-100) mm 改为 (0,0,-60) mm
- **几何结构**：两个盒体的尺寸将更新为您指定的数值（A: 40x40x40 mm³, B: 20x20x20 mm³）

如果您确认应用这些修改，请直接回复“确认”；如果您希望保留原来的值，请回复“保持原值”。
\item User: 确认
\item Assistant: 配置已完成。当前已确认的设置包括：源位置、布尔 solid B 尺寸、布尔 solid A 尺寸和输出路径。
\end{itemize}
\textbf{Internal Chain (Natural Language):}
\begin{itemize}
\item Turn 1: action=summarize\_progress, phase=geometry, complete=False.
\item Missing fields: geometry.params.bool\_a\_x, geometry.params.bool\_a\_y, geometry.params.bool\_a\_z, geometry.params.bool\_b\_x, geometry.params.bool\_b\_y, geometry.params.bool\_b\_z, materials.volume\_material\_map, source.particle, source.energy, physics.physics\_list, output.format, output.path, geometry.ask.boolean.solid\_a\_size, geometry.ask.boolean.solid\_b\_size.
\item LLM used=True; fallback=none; backend=llm\_slot\_frame+runtime\_semantic.
\item Turn 2: action=confirm\_overwrite, phase=materials, complete=False.
\item Missing fields: none.
\item LLM used=True; fallback=none; backend=llm\_slot\_frame+runtime\_semantic.
\item Turn 3: action=finalize, phase=source, complete=True.
\item Missing fields: none.
\item LLM used=True; fallback=none; backend=llm\_semantic\_frame+runtime\_semantic+confirmed\_pending\_overwrite.
\end{itemize}
\subsection{Case MTM19 (ring\_complete)}
\textbf{Dialogue:}
\begin{itemize}
\item User: Create a complete ring setup: 16 modules, each 6 mm x 6 mm x 2 mm, radius 32 mm, clearance 0.5 mm, material G4\_Si; point gamma source 0.511 MeV at (0,0,0) mm; physics QBBC; output json.
\item Assistant: Your configuration is complete. I have confirmed the volume-to-material binding, output path, clearance, and radius.
\end{itemize}
\textbf{Internal Chain (Natural Language):}
\begin{itemize}
\item Turn 1: action=finalize, phase=materials, complete=True.
\item Missing fields: none.
\item LLM used=True; fallback=none; backend=llm\_slot\_frame+runtime\_semantic.
\end{itemize}
\subsection{Case MTM20 (grid\_complete)}
\textbf{Dialogue:}
\begin{itemize}
\item User: 请创建完整的阵列配置：3 x 3 盒体阵列，每个模块 12 mm x 12 mm x 3 mm，pitch\_x 15 mm，pitch\_y 15 mm，间隙 1 mm，材料 G4\_CsI；点源 gamma 0.662 MeV，位置 (0,0,-90) mm，方向 +z；物理列表 QBBC；输出 root。
\item Assistant: 配置已完成。当前已确认：体积与材料绑定、输出路径、clearance、y 方向 pitch。
\end{itemize}
\textbf{Internal Chain (Natural Language):}
\begin{itemize}
\item Turn 1: action=finalize, phase=materials, complete=True.
\item Missing fields: none.
\item LLM used=True; fallback=none; backend=llm\_slot\_frame+runtime\_semantic.
\end{itemize}
\subsection{Case MTM21 (layered\_detector)}
\textbf{Dialogue:}
\begin{itemize}
\item User: I want a stacked detector slab, but the individual layer thicknesses are not fixed yet.
\item Assistant: Here's what we have so far:

**Geometry:** We've set up a stacked detector slab structure. We still need the footprint (size) of each layer and their individual thicknesses.

**Materials:** We still need to assign materials to the detector volumes.

**Source:** We still need the source type (like a beam or a point) and the type of particle to simulate.

**Physics \& Output:** We still need to choose the physics models and specify the output format and file location.
\item User: Use stack footprint 25 mm x 25 mm, thicknesses 1 mm, 2 mm, 1 mm, layer clearance 0.2 mm, outer box 30 mm x 30 mm x 12 mm, container clearance 0.5 mm, material G4\_Si; beam electron 5 MeV from (0,0,-50) mm along +z; physics FTFP\_BERT; output json.
\item Assistant: This would overwrite a geometry you already confirmed.

**Current geometry:**
*   Stack footprint: \textasciitilde{}16.3 mm x \textasciitilde{}10.3 mm
*   Layer thicknesses: \textasciitilde{}6.68 mm, \textasciitilde{}1.19 mm, \textasciitilde{}5.59 mm
*   Layer clearance: \textasciitilde{}0.73 mm
*   Outer box: \textasciitilde{}18.1 mm x \textasciitilde{}45.5 mm x \textasciitilde{}16.7 mm
*   Container clearance: \textasciitilde{}0.87 mm

**New geometry from your latest request:**
*   Stack footprint: 25.0 mm x 25.0 mm
*   Layer thicknesses: 1.0 mm, 2.0 mm, 1.0 mm
*   Layer clearance: 0.2 mm
*   Outer box: 30.0 mm x 30.0 mm x 12.0 mm
*   Container clearance: 0.5 mm

Please reply **'confirm'** to apply this change, or **'keep existing'** to keep the current values.
\item User: confirm
\item Assistant: Configuration complete. The following geometry items are now confirmed: stack container clearance, stack container size, stack layer clearance, and stack layer thicknesses.
\end{itemize}
\textbf{Internal Chain (Natural Language):}
\begin{itemize}
\item Turn 1: action=summarize\_progress, phase=materials, complete=False.
\item Missing fields: materials.volume\_material\_map, source.type, source.particle, source.energy, source.position, source.direction, physics.physics\_list, output.format, output.path, geometry.ask.stack.footprint, geometry.ask.stack.thicknesses, geometry.ask.stack.layer\_clearance, geometry.ask.stack.parent\_size, geometry.ask.stack.nest\_clearance.
\item LLM used=True; fallback=none; backend=llm\_semantic\_frame+runtime\_semantic.
\item Turn 2: action=confirm\_overwrite, phase=source, complete=False.
\item Missing fields: none.
\item LLM used=True; fallback=none; backend=llm\_slot\_frame+runtime\_semantic.
\item Turn 3: action=finalize, phase=physics, complete=True.
\item Missing fields: none.
\item LLM used=True; fallback=none; backend=llm\_semantic\_frame+runtime\_semantic+confirmed\_pending\_overwrite.
\end{itemize}
\subsection{Case MTM22 (embedded\_sensor)}
\textbf{Dialogue:}
\begin{itemize}
\item User: 请创建完整的嵌套配置：外盒 80 mm x 80 mm x 80 mm，内嵌铅圆柱半径 15 mm、半长 25 mm，嵌套间隙 1 mm；材料 G4\_Pb；点源 gamma 1 MeV，位置 (0,0,-100) mm，方向 +z；物理列表 Shielding；输出 csv。
\item Assistant: 配置已完成！当前已确认的设置包括：体积与材料绑定、输出路径、嵌套间隙（clearance）以及圆柱半径。
\end{itemize}
\textbf{Internal Chain (Natural Language):}
\begin{itemize}
\item Turn 1: action=finalize, phase=materials, complete=True.
\item Missing fields: none.
\item LLM used=True; fallback=none; backend=llm\_slot\_frame+runtime\_semantic.
\end{itemize}
\subsection{Case MTM23 (layered\_shield\_complete)}
\textbf{Dialogue:}
\begin{itemize}
\item User: Create a complete shell configuration: inner radius 12 mm, thicknesses 4 mm and 6 mm, half length 35 mm, nested core radius 8 mm, nested core half length 25 mm, clearance 1 mm, material G4\_Fe; isotropic gamma source 0.8 MeV at (0,0,0) mm; physics QBBC; output hdf5.
\item Assistant: Your configuration is complete. The following items have been confirmed: box dimensions (x, y, z) and the volume-to-material binding.
\end{itemize}
\textbf{Internal Chain (Natural Language):}
\begin{itemize}
\item Turn 1: action=finalize, phase=materials, complete=True.
\item Missing fields: none.
\item LLM used=True; fallback=none; backend=llm\_slot\_frame+runtime\_semantic.
\end{itemize}
\subsection{Case MTM24 (boolean\_complete)}
\textbf{Dialogue:}
\begin{itemize}
\item User: 请创建完整布尔配置：做并集布尔体，A 为 30 mm x 30 mm x 20 mm 盒体，B 为 10 mm x 10 mm x 40 mm 盒体，与 A 做 union；材料 G4\_Al；束流 gamma 0.5 MeV，位置 (0,0,-70) mm，方向 +z；物理列表 QBBC；输出 json。
\item Assistant: 配置已完成。当前已确认的参数包括：盒体在 x、y、z 方向的尺寸，以及体积与材料的绑定关系。
\end{itemize}
\textbf{Internal Chain (Natural Language):}
\begin{itemize}
\item Turn 1: action=finalize, phase=materials, complete=True.
\item Missing fields: none.
\item LLM used=True; fallback=none; backend=llm\_slot\_frame+runtime\_semantic.
\end{itemize}
\end{document}